\section{Introduction}

A spectrograph disperses the light of a source in wavelength, giving access to information that an image does not carry: chemical composition, temperature, radial velocity, or whether the emission is a continuum or a set of lines. The basic layout is a slit, a collimator, a dispersing element and a camera that focuses each wavelength at a different position of the detector.

The instrument assembled in this experiment is BESOS (\emph{BEst Simple Optical Spectrograph}), a low-resolution design by the CAOS group~\cite{besos}. Its dispersing element is an equilateral SF18 prism, so the dispersion comes from the variation of the refractive index with wavelength. Compared with a grating, a prism gives lower resolution and a strongly non-linear dispersion (larger in the blue than in the red), but it sends all the light into a single spectrum.

The goal is to assemble the instrument and to measure in the laboratory its \textbf{spectral range}, \textbf{dispersion} and \textbf{spectral resolution}, comparing them with the first-order expectation~\cite{guide}.
